\begin{center}
\begin{figure}[H]
\centering
\begin{tikzpicture}[
  cell/.style={rectangle, draw, minimum size=7mm, anchor=center},
  node/.style={rectangle, draw, fill=blue!10, minimum height=6mm, minimum width=12mm, font=\small},
  thick, >=stealth
]
  % Tabulka
  \matrix (m) [matrix of math nodes, nodes={draw, minimum size=7mm, anchor=center}, row sep=-\pgflinewidth] {
    0 \\ 1 \\ 2 \\ 3 \\ 4 \\
  };

  % Seznamy
  \node[node, right=12pt of m-2-1] (n15) {15};
  \draw[->] (m-2-1) -- (n15);
  \node[node, right=10pt of n15] (n5) {5};
  \draw[->] (n15) -- (n5);

  \node[node, right=12pt of m-4-1] (n8) {8};
  \draw[->] (m-4-1) -- (n8);

  \node[node, right=12pt of m-5-1] (n14) {14};
  \draw[->] (m-5-1) -- (n14);
  \node[node, right=10pt of n14] (n4) {4};
  \draw[->] (n14) -- (n4);

\end{tikzpicture}
\caption{Schéma hešování s řetězením. Přihrádka 1 obsahuje spojový seznam prvků $[15 \to 5]$, přihrádky 0 a 2 jsou prázdné.}
\end{figure}
\end{center}